% Compile with pdfLaTeX or Tectonic. No external figures or bibliography required.
% Numerical sources: design_summary.csv, verification_summary.csv,
% simulation_log.txt, and calculate_design.m in this directory.
% All displayed values are rounded from the saved full-precision results.
\documentclass[11pt,a4paper]{article}
\usepackage[T1]{fontenc}
\usepackage{lmodern}
\usepackage[margin=20mm]{geometry}
\usepackage{amsmath,amssymb,booktabs}
\usepackage{microtype}
\setlength{\parindent}{0pt}
\setlength{\parskip}{4pt}
\setlength{\abovedisplayskip}{7pt}
\setlength{\belowdisplayskip}{7pt}
\newcommand{\step}[1]{\par\vspace{8pt}{\large\bfseries #1}\par\vspace{2pt}}
\newcommand{\ohm}{\,\Omega}
\newcommand{\nH}{\,\mathrm{nH}}
\newcommand{\nF}{\,\mathrm{nF}}
\newcommand{\radsec}{\,\mathrm{rad\,s^{-1}}}
\pagestyle{plain}

\begin{document}
{\LARGE\bfseries LCC-C Inverse Design}\par
{\large Step-by-Step Calculation and Simulation Results}\par
\vspace{3pt}
The calculation follows the sequence
\[
 (f_0,R_L,L_1,L_2,Q_1,Q_2)
 \;\longrightarrow\; (R_1,R_2)
 \;\longrightarrow\; k
 \;\longrightarrow\; M
 \;\longrightarrow\; (L,C).
\]
The target input impedance is $Z_{\mathrm{in}}=30+j0\ohm$ at $f_0$.
Here, $L_1$ and $L_2$ are coil self-inductances; $L$ is the additional
transmitter compensation inductance.

\step{1. Known parameters and angular frequency}
\begin{center}
\begin{tabular}{@{}ll@{}}
\toprule
Quantity & Value \\
\midrule
Operating frequency, $f_0$ & $13.56\,\mathrm{MHz}$ \\
Load resistance, $R_L$ & $12.5\ohm$ \\
Coil self-inductances, $L_1=L_2$ & $486.02\nH$ \\
Unloaded coil quality factors & $Q_1=951.1$, $Q_2=931$ \\
Input impedance requirement & $30+j0\ohm$ \\
\bottomrule
\end{tabular}
\end{center}
\[
 \omega_0=2\pi f_0=2\pi(13.56\times10^6)
          =8.5199992765\times10^7\radsec.
\]

\step{2. Winding resistances from the unloaded Q values}
Using $Q_i=\omega_0L_i/R_i$ at the operating frequency,
\begin{align*}
 R_1&=\frac{\omega_0L_1}{Q_1}
     =\frac{(8.5199992765\times10^7)(486.02\times10^{-9})}{951.1}
     =0.043537903989\ohm,\\[3pt]
 R_2&=\frac{\omega_0L_2}{Q_2}
     =\frac{(8.5199992765\times10^7)(486.02\times10^{-9})}{931}
     =0.044477873774\ohm.
\end{align*}

\step{3. Coupling coefficient k from the optimum-load condition}
For the series-tuned receiver, choose the given load to be the
efficiency-optimum load:
\[
 R_L=R_{\mathrm{opt}}=R_2\sqrt{1+k^2Q_1Q_2}.
\]
This condition is required to determine $k$; the Q values alone do not fix it.
\begin{align*}
 k&=\sqrt{\frac{(R_L/R_2)^2-1}{Q_1Q_2}}\\
  &=\sqrt{\frac{(12.5/0.044477873774)^2-1}{951.1\times931}}
   =\boxed{0.298658799019}.
\end{align*}
Substitution into the optimum-load equation gives
$R_{\mathrm{opt}}=12.500000000\ohm$.

\step{4. Mutual inductance M}
\begin{align*}
 M&=k\sqrt{L_1L_2}\\
  &=0.298658799019\sqrt{(486.02\times10^{-9})(486.02\times10^{-9})}\,\mathrm{H}\\
  &=\boxed{145.154149499\nH}.
\end{align*}

\newpage
\step{5. Reflected resistance and compensation inductance L}
At receiver series resonance, the reflected resistance and total
transmitter-branch resistance are
\begin{align*}
 R_{\mathrm{ref}}
   &=\frac{(\omega_0M)^2}{R_2+R_L}
     =\frac{(12.3671324872)^2}{0.044477873774+12.5}
     =12.1922942903\ohm,\\
 R_{\mathrm{branch}}
   &=R_1+R_{\mathrm{ref}}
     =0.043537903989+12.1922942903
     =12.2358321943\ohm.
\end{align*}
Under the LCC tuning conditions
$\omega_0^2LC=1$ and $C_1=1/[\omega_0^2(L_1-L)]$,
the input impedance is real:
\[
 Z_{\mathrm{in}}=\frac{(\omega_0L)^2}{R_{\mathrm{branch}}}=30\ohm.
\]
Therefore,
\[
 L=\frac{\sqrt{30R_{\mathrm{branch}}}}{\omega_0}
  =\frac{\sqrt{30\times12.2358321943}}{8.5199992765\times10^7}
  =\boxed{224.873265022\nH}.
\]

\step{6. Compensation capacitances C, C1 and C2}
The shunt compensation capacitor follows directly from $L$:
\[
 C=\frac{1}{\omega_0^2L}
  =\frac{1}{(8.5199992765\times10^7)^2(224.873265022\times10^{-9})}
  =\boxed{0.612608556712\nF}.
\]
The two series capacitors complete the LCC-C tuning:
\begin{align*}
 C_1&=\frac{1}{\omega_0^2(L_1-L)}
     =\frac{1}{\omega_0^2(261.146734978\times10^{-9})}
     =0.527516786069\nF,\\
 C_2&=\frac{1}{\omega_0^2L_2}
     =\frac{1}{\omega_0^2(486.02\times10^{-9})}
     =0.283443657315\nF.
\end{align*}

\step{7. Simulation verification of the calculated design}
The saved Simulink load sweep peaks at $R_L=12.5\ohm$.
At $13.56\,\mathrm{MHz}$, the frequency-domain model gives
\[
 Z_{\mathrm{in}}=30-j\,2.4869\times10^{-14}\ohm\approx30+j0\ohm,
 \qquad \eta=99.290877241\%.
\]
The independent Simscape model uses a $1\,\mathrm{V_{rms}}$ source,
the \texttt{ode23t} solver and maximum step $T/25$.
At $40\,\mu\mathrm{s}$, it gives
\begin{align*}
 Z_{\mathrm{in}}&=29.992296506-j\,0.001738268\ohm,\\
 P_{\mathrm{in}}&=0.033340115\,\mathrm{W},\qquad
 P_{\mathrm{load}}=0.033098124\,\mathrm{W},\qquad
 \eta=99.274173\%.
\end{align*}
\begin{center}
\small
\begin{tabular}{@{}llll@{}}
\toprule
Check & Result & Requirement & Status \\
\midrule
Real input-impedance error & $0.0256783\%$ & $<1\%$ & PASS \\
Absolute input reactance & $0.001738268\ohm$ & $<0.3\ohm$ & PASS \\
Adjacent 20-cycle phasor change & $0.0179217\%$ & $<0.1\%$ & PASS \\
Input and load powers & Finite and positive & Both $>0$ & PASS \\
\bottomrule
\end{tabular}
\end{center}
{\small Model scope: ideal compensation components with winding losses
$R_1$ and $R_2$; switching and rectification losses are excluded.
$M$ is an inferred design target, not a measured coil-pair value.}

\vspace{3pt}
{\footnotesize Reproduce in MATLAB R2026b using \texttt{run\_serial\_link}.\par
Numerical records: \texttt{design\_summary.csv},
\texttt{verification\_summary.csv} and \texttt{simulation\_log.txt}.}
\end{document}
